\documentclass[11pt]{article}
\usepackage[utf8]{inputenc}
\usepackage[T1]{fontenc}
\title{Organización de la interfaz de TerraPulse AI}
\author{TerraPulse AI}
\date{}
\begin{document}
\maketitle
\section{Objetivo}
Separar las acciones del mapa de la lectura de resultados, sin cambiar los cálculos SAR o solares. El trabajo se originó en el worktree \texttt{data/worktrees/frontend-workspace-ui} y fue integrado por archivo al directorio principal, con respaldo en \texttt{data/integration-backups/frontend-c9a646629763}. No se realizó commit ni push.
\section{Uso}
Las acciones Copia con Calor, Agua DG2 y Sismo permanecen en la barra contextual y dejan de repetirse en las pestañas del mapa. El panel solar solicita una ROI confirmada, fecha, horario ilustrativo, zona horaria IANA explícita e intensidad física. Sin zona válida, la ejecución queda deshabilitada. La opacidad controla solo la presentación de la capa. Al cambiar ROI, fecha, zona, horario o intensidad, se invalida la vista de una corrida anterior.
El botón \textit{Ver análisis completo} abre un diálogo amplio con secciones de resumen, distribución de irradiancia, gráficas disponibles, capas y muestra geoespacial, metadatos y exportaciones. Las gráficas NISAR, cuando existen, se identifican como producto distinto; no se inventan series solares. Se puede cerrar con Escape y volver al mapa sin perder la configuración.
\section{Implementación}
\texttt{Header.tsx} conserva las acciones de copia; \texttt{MapCanvas.tsx} usa la misma condición de visibilidad SAR para capa y leyenda, y retira la escala solar porcentual que confundía intensidad física con W/m\textsuperscript{2}. \texttt{RightPanel.tsx} muestra la leyenda solar y los datos numéricos recibidos de los endpoints de análisis y muestra del GeoTIFF, y reutiliza la descarga GeoTIFF y la vista PNG. No obtiene medidas del color PNG. \texttt{SolarAnalysis.css} define el diálogo adaptable; \texttt{index.css} permite ajustar barras, gráficas y leyendas. La resolución y el tamaño de píxel se redondean solo en pantalla; los datos fuente no cambian. El explorador identifica las pilas de Chiapas como caso de referencia, no como ROI activa.
\section{Limitaciones}
No existe servicio que deduzca la zona horaria de una ROI: el usuario debe seleccionar una IANA explícita, incluso para ROIs multizona. La validación del navegador no garantiza que la hora elegida exista durante cambios de horario de verano; el backend puede rechazarla. La definición de noche en corridas históricas puede diferir de las nuevas; revisar \texttt{code\_version} y elevación geométrica antes de comparar. El análisis disponible contiene una capa directa incidente, sin series solares adicionales. Este trabajo no verifica la respuesta de servicios remotos ni integra los cambios pendientes del backend de otra sesión.
\section{Verificación}
En el worktree aislado se observó RED en pruebas de duplicación, secciones de análisis y presentación antes de corregir esos casos; la verificación independiente no presenció esos RED y las pruebas de overlays son principalmente estructurales. Después de integrar en principal, \texttt{node --test tests/*.test.mjs} aprobó 103 pruebas frontend, \texttt{npm run build} compiló, y \texttt{PYTHONPATH=backend:.} con el Python del entorno \texttt{nisar} y \texttt{pytest -q pipeline/tests/test\_solar\_*.py} aprobó 230 pruebas. Un puente de contrato sin red, en \texttt{data/integration-backups/frontend-c9a646629763/contract\_bridge.py} y \texttt{contract\_bridge.mjs}, verificó respuestas reales del cliente HTTP Django (POST, análisis, muestra, errores 400 y 404) consumidas por \texttt{web/src/api/client.ts}; utiliza una superficie de prueba simulada, no NASADEM real. Una comparación diferencial de cinco funciones públicas no solares contra el cliente previo no encontró cambios en rutas o errores de respuesta ordinarios. Las descargas y el navegador permanecen sin comprobación funcional. La interfaz presenta los errores 400/404 de análisis y muestra y limpia el error de POST al editar parámetros; los enlaces nativos de exportación TIFF y vista PNG todavía no muestran dentro de la interfaz el estado HTTP 404 del archivo descargado. \texttt{npm run build} compiló con avisos sobre configuración Vite y tamaño del paquete. \texttt{git diff --check} no reportó errores. \texttt{npm run lint} terminó con cuatro errores en archivos ajenos a esta modificación (\texttt{ScientificFiguresViewer.tsx} y \texttt{utils/seismic.ts}) y advertencias de dependencias de hooks. No hubo navegador disponible para comprobar consola, teclado, exportaciones reales y tamaños de viewport: esa verificación queda pendiente y las pruebas de código fuente no la reemplazan. No se realizó commit ni push.
\end{document}
